\section{The intrinsic algebra of the original normal filtration}

Write $R=\mathcal O_{\mathbb C^{d+t},0}$ and
$B=\mathcal O_{\mathbb C^d,0}$, with the actual tangent pullback giving
the $B$-algebra structure on $R$. Let $J$ be the original normal ideal and
let $I\subset R$ be any original whole germ ideal. Its existing literal
degree-$q$ module is the image of $J^q$ in $R/(I+J^{q+1})$. The algebra
below is constructed from multiplication of these genuine representatives;
the polynomial quotient is a subsequent theorem about that algebra. This
commutative-algebra construction is the explicit intrinsic algebra adaptation
needed by the relative normal-flatness and Hilbert--Samuel route; the quotient
never defines the multiplication.

The source-facing objective is the whole local-ring Hilbert--Samuel
quotient used in Bierstone--Milman, ``Use of the Hilbert--Samuel function''.
The following constructions supply its relative normal-filtration algebra
ingredient directly from the original germ ideals. They add the intrinsic
multiplication to the preceding literal-module dictionary; they do not assert
the source's semicontinuity, admissible-centre theorem, or descent under
blowing up.

\begin{definition}[Actual multiplication of original normal layers]
\label{hir:relative-normal-intrinsic-layer-multiplication}
\source{bierstone-milman08-functoriality}
\lean{Hironaka.HironakaRelativeNormalGradedAlgebra.originalNormalGradedPieceMul}
\leanok
\uses{hir:relative-normal-graded-literal-layers}
For $x$ in degree $p$ and $y$ in degree $q$, choose actual representatives
$g\in J^p$ and $h\in J^q$ of their whole quotient classes and define
$x\cdot y$ to be the class of $gh$ in the image of $J^{p+q}$ in
$R/(I+J^{p+q+1})$.
\end{definition}
\begin{proof}
The representatives exist by the literal image-module definition.
Their product lies in $J^pJ^q=J^{p+q}$, hence gives an element of the
stated actual layer. Independence from these representatives is proved in
the following lemma.
\end{proof}

\begin{lemma}[The whole original remainder makes multiplication well-defined]
\label{hir:relative-normal-intrinsic-layer-product-independent}
\source{bierstone-milman08-functoriality}
\lean{Hironaka.HironakaRelativeNormalGradedAlgebra.originalNormalGradedPieceMul_of_representatives}
\leanok
\uses{hir:relative-normal-intrinsic-layer-multiplication}
The actual product class is independent of the chosen whole original
representatives. In particular, the product of the degree-$p$ class of
$g\in J^p$ and the degree-$q$ class of $h\in J^q$ is exactly the
degree-$(p+q)$ class of their actual product $gh$.
\end{lemma}
\begin{proof}
Suppose $g-g'\in I+J^{p+1}$ and $h-h'\in I+J^{q+1}$, with each
representative in its asserted original power. For $v\in I+J^{p+1}$
and $w\in J^q$, write $v=i+u$ with $i\in I$ and
$u\in J^{p+1}$. Then $iw\in I$ and
$uw\in J^{p+1}J^q=J^{p+q+1}$, so
$vw\in I+J^{p+q+1}$. Apply this first to $(g-g')h$, and then,
with the two degrees interchanged, to $g'(h-h')$. Their sum is
$gh-g'h'$, which therefore belongs to the whole next-power remainder
$I+J^{p+q+1}$. The two product classes coincide in exactly the original
quotient that defines the layer.
\end{proof}

\begin{theorem}[The literal original direct sum is an intrinsic graded algebra]
\label{hir:relative-normal-intrinsic-graded-algebra}
\source{bierstone-milman08-functoriality}
\lean{Hironaka.HironakaRelativeNormalGradedAlgebra.originalNormalAssociatedGradedCommRing,Hironaka.HironakaRelativeNormalGradedAlgebra.originalNormalAssociatedGradedAlgebra,Hironaka.HironakaRelativeNormalGradedAlgebra.originalNormalAssociatedGraded_lof_mul_lof,Hironaka.HironakaRelativeNormalGradedAlgebra.originalNormalAssociatedGraded_algebraMap}
\leanok
\uses{hir:relative-normal-intrinsic-layer-product-independent}
The existing literal direct sum of all original normal layers is a
commutative $B$-algebra, with the product of its degree-$p$ and
degree-$q$ summands given by the preceding actual multiplication and lying
in degree $p+q$. Its unit is the actual class of $1$ in degree zero.
Its $B$-scalar action is the existing tangent-pullback action on each layer.
\end{theorem}
\begin{proof}
The family of actual powers $J^q\subset R$, with ordinary multiplication
and degree-zero unit $1$, is a graded commutative ring. The true quotient
map from each $J^q$ onto its literal layer is additive and $B$-linear,
and the preceding lemma makes these maps multiplicative for the already
defined actual layer products. Every layer class has an actual power
representative. Thus associativity, commutativity, the two unit laws, and
both distributive laws follow by lifting the classes to these genuine
power representatives and applying the corresponding identities in $R$.
The degree-zero scalar class is the class of the actual tangent pullback of
the scalar; multiplying it by a degree-$q$ representative gives exactly
the existing scalar action on the whole quotient. Extend the homogeneous
products to finite sums by distributivity. All sums involved are finite,
and these proved homogeneous laws give the stated commutative algebra on
the original direct sum.
\end{proof}

\begin{lemma}[The actual entire polynomial symbol map is an algebra homomorphism]
\label{hir:relative-normal-intrinsic-polynomial-symbol-algebra-hom}
\source{bierstone-milman08-functoriality}
\lean{Hironaka.HironakaRelativeNormalGradedAlgebra.normalAssociatedSymbolAlgHom,Hironaka.HironakaRelativeNormalGradedAlgebra.normalAssociatedSymbolAlgHom_toLinearMap,Hironaka.HironakaRelativeNormalGradedAlgebra.normalAssociatedSymbolAlgHom_ker}
\leanok
\uses{hir:relative-normal-intrinsic-graded-algebra,hir:relative-normal-graded-homogeneous-surjection,hir:relative-normal-graded-whole-initial-ideal}
The existing true polynomial symbol map from $B[Y_1,\ldots,Y_t]$ to the
literal original associated module is a unital $B$-algebra homomorphism
for its intrinsic multiplication. Its complete kernel is the existing
whole original normal initial ideal.
\end{lemma}
\begin{proof}
For homogeneous $P$ of degree $p$ and $Q$ of degree $q$, their symbol
classes are the actual classes of their genuine normal evaluations in
their respective layers. The intrinsic product lemma identifies their
product with the class of the product of these evaluations in degree
$p+q$. Actual normal polynomial evaluation is an algebra homomorphism,
so this is exactly the symbol of the homogeneous polynomial $PQ$.
Write arbitrary $P$ and $Q$ as their finite sums of monomials, each of
which is homogeneous of its actual exponent degree. Induct through these
two finite sums. Additivity of the existing symbol map and the two
distributive laws reduce every induction step to the already proved
homogeneous product identity. This gives the same multiplicative identity
for the two whole polynomials.
The symbol of $1$ is its actual degree-zero class, which is the
intrinsic unit. A constant coefficient has the actual degree-zero tangent
pullback class, proving preservation of scalars. The underlying map is
the existing entire symbol map, so its previously proved complete
kernel dictionary gives exactly the existing original initial ideal.
\end{proof}

\begin{theorem}[The polynomial quotient is the intrinsic original associated algebra]
\label{hir:relative-normal-intrinsic-associated-quotient-algebra-dictionary}
\source{bierstone-milman08-functoriality}
\lean{Hironaka.HironakaRelativeNormalGradedAlgebra.originalNormalAssociatedGradedQuotientAlgEquiv,Hironaka.HironakaRelativeNormalGradedAlgebra.originalNormalAssociatedGradedQuotientAlgEquiv_mk,Hironaka.HironakaRelativeNormalGradedAlgebra.originalNormalAssociatedGradedQuotientAlgEquiv_toLinearEquiv,Hironaka.HironakaRelativeNormalGradedAlgebra.originalNormalAssociatedGradedQuotientAlgEquiv_homogeneous}
\leanok
\uses{hir:relative-normal-intrinsic-polynomial-symbol-algebra-hom,hir:relative-normal-graded-internal-finite-relations}
There is a $B$-algebra isomorphism from the whole polynomial quotient by
the existing original normal initial ideal onto the intrinsic original
associated algebra. It sends every actual polynomial numerator to its
same existing entire original symbol, retaining every homogeneous degree.
\end{theorem}
\begin{proof}
The original homogeneous symbol surjections give surjectivity of the
entire symbol map onto the literal direct sum. Its intrinsic algebra
homomorphism has exactly the existing whole original initial ideal as
kernel. The algebraic first isomorphism theorem therefore gives the
asserted quotient isomorphism. Its defining quotient map sends every
polynomial numerator to that same actual symbol. For a homogeneous
numerator the original component dictionary places its symbol in its
actual same-degree summand, proving the degree assertion. Since every original quotient class has a polynomial numerator, the numerator identity also proves that the underlying linear isomorphism is exactly the previous literal quotient-module dictionary. The ring and
scalar structures used here were constructed before this quotient
comparison from genuine original representatives.
\end{proof}
